\documentclass[11pt]{article}
\usepackage[margin=1in]{geometry}
\usepackage{amsmath,amssymb}
\usepackage[hidelinks]{hyperref}

\newcommand{\D}{\mathbb D}
\newcommand{\T}{\mathbb T}
\newcommand{\Po}{\mathcal P}
\newcommand{\BC}{\mathrm{BC}}

\title{A Literature-Implied Answer to the Kriete--MacCluer Beurling--Carleson Regime}
\author{Run fa\_banach\_001}
\date{2026-07-18}

\begin{document}
\maketitle

\paragraph{Status.}
This is a lightweight literature-implied status note:
\[
\texttt{literature\_implied\_answer.}
\]
It records that the future-work direction flagged in Malman,
\emph{Revisiting mean-square approximation by polynomials in the unit disk},
arXiv:2304.01400, is answered in the principal standard-weighted Bergman
setting, and in the explicit subexponential gauge model, by Malman's later
paper \emph{Weighted Korenblum-Roberts theory}, arXiv:2503.20054.

\paragraph{Source question.}
In arXiv:2304.01400, Introduction, subsection ``Work of Kriete and
MacCluer'', the source explains that Kriete and MacCluer's Conjecture 2 is
settled by the paper in the fast-decay regime
\[
\liminf_{x\to 0^+} x\log(1/G(x))>0.
\]
The next paragraph records their Conjecture 1 for the limiting-zero regime.
The expected rule is analogous to the fast-decay theorem, but with intervals
replaced by Beurling--Carleson sets, and the source says that a conclusive
answer is future work.

\paragraph{Supporting theorem.}
In arXiv:2503.20054, the supporting paper defines the associated
Beurling--Carleson family
\[
\operatorname{Assoc}(w)
 =\left\{E\in \BC:\int_E \log w\,dm>-\infty\right\},
\]
then defines the core of \(w\) as the maximal-measure union generated by such
sets and the residual part as the carrier of \(w\) outside the core.  Its
Theorem~\(\mathrm{T{:}ThomsonDecompTheorem}\) states that, for
\[
d\mu=dA_\alpha+w\,dm,\qquad \alpha>-1,
\]
and every \(0<t<\infty\),
\[
\Po^t(\mu)=
\Po^t(dA_\alpha+w_c\,dm)\oplus L^t(w_r\,dm),
\]
where \(w_c\) and \(w_r\) are the core and residual restrictions of \(w\), and
the first summand is irreducible.  Thus full splitting off the boundary
Lebesgue part occurs exactly when the Beurling--Carleson core contribution
vanishes, and in general the theorem gives the promised Beurling--Carleson
replacement for the interval residual decomposition from the fast-decay
regime.

\paragraph{Gauge version and match to the 2023 discussion.}
The same supporting paper also formulates the generalized measure
\[
d\mu=G_{ah}\,dA+w\,dm,\qquad
G_{ah}(z)=\exp\!\left(-a\,\frac{h(1-|z|)}{1-|z|}\right),
\]
with the associated \(h\)-Beurling--Carleson family.  It states explicitly
that for \(h(x)=x^\beta\), \(0<\beta<1\), one has
\[
\Po^t(\mu)=
\Po^t(G_{ah}\,dA+w_c\,dm)\oplus L^t(w_r\,dm),
\]
again with an irreducible first summand.  This includes
\[
G_{ah}(z)=\exp\!\left(-\frac{a}{(1-|z|)^{1-\beta}}\right),
\]
which is the subexponential model discussed in the sharpness section of
arXiv:2304.01400.

\paragraph{Why this is literature-implied rather than already-answered.}
The supporting paper cites arXiv:2304.01400 and explains the relation between
the fast interval regime and the slower associated-Beurling--Carleson regime,
but it does not label its theorem as ``Kriete--MacCluer Conjecture 1''.  The
identification is therefore a direct reformulation by this run: matching the
source's requested replacement of intervals by Beurling--Carleson sets with
the later core/residual decomposition theorem.

\paragraph{Scope limitations.}
This packet should not be used as a claim that every radial disk weight \(G\)
with
\[
\lim_{x\to0^+}x\log(1/G(x))=0
\]
is covered.  The supporting theorem covers the standard weighted Bergman area
case and the stated gauge families, including \(h(x)=x^\beta\).  Any arbitrary
radial-weight extension would still require checking the regularity/gauge
hypotheses or a separate theorem.

\paragraph{References.}
\begin{itemize}
\item B. Malman, \emph{Revisiting mean-square approximation by polynomials in
the unit disk}, arXiv:2304.01400.
\item B. Malman, \emph{Weighted Korenblum-Roberts theory}, arXiv:2503.20054.
\item T. L. Kriete and B. D. MacCluer, \emph{Mean-square approximation by
polynomials on the unit disk}, Transactions of the American Mathematical
Society 322 (1990), 1--34.
\end{itemize}

\end{document}
